% =====================================================================
%  ISR-401 Ingenieria de Requerimientos - UTEQ
%  TA-PFC-E1: vision, alcance, partes interesadas, contexto y plan de
%  elicitacion
%  Compilar: pdflatex main -> bibtex main -> pdflatex main -> pdflatex main
%  (o simplemente: latexmk -pdf main.tex)
% =====================================================================
\documentclass[11pt,a4paper]{article}

% ---------- Codificacion e idioma ----------
\usepackage[utf8]{inputenc}
\usepackage[T1]{fontenc}
\usepackage{lmodern}
\usepackage[spanish,es-tabla,es-noshorthands]{babel}

% ---------- Pagina ----------
\usepackage[margin=2.5cm]{geometry}
\usepackage{fancyhdr}
\usepackage{setspace}
\onehalfspacing

% ---------- Tablas, listas, figuras ----------
\usepackage{graphicx}
\graphicspath{{figuras/}}
\usepackage{booktabs}
\usepackage{tabularx}
\usepackage{longtable}
\usepackage{array}
\usepackage{multirow}
\usepackage[table]{xcolor}
\usepackage{enumitem}
\usepackage{caption}
\usepackage{float}
\usepackage{amssymb} % \square en los instrumentos

% ---------- Diagramas (fuentes en modelos/) ----------
\usepackage{tikz}
\usetikzlibrary{positioning,arrows.meta,shapes.geometric,fit,calc}

% ---------- Enlaces ----------
\usepackage[hidelinks]{hyperref}
\usepackage{url}

% ---------- Metadatos del PDF ----------
\hypersetup{
  pdftitle={TA-PFC-E1: Vision, alcance, partes interesadas, contexto y plan de elicitacion},
  pdfauthor={Equipo ISR401-ContX},
  pdfsubject={Ingenieria de Requerimientos ISR-401, UTEQ},
  pdfkeywords={requisitos, elicitacion, partes interesadas, ISO/IEC/IEEE 29148}
}

% ---------- Colores y macros ----------
\definecolor{uteqverde}{RGB}{16,122,62}
\definecolor{pendiente}{RGB}{200,0,0}
% \pendiente{...} marca lo que el equipo debe completar.
% Antes de entregar, busque "\pendiente" en todo el proyecto: no debe quedar ninguno.
\newcommand{\pendiente}[1]{\textcolor{pendiente}{\textbf{[PENDIENTE: #1]}}}
\newcolumntype{L}[1]{>{\raggedright\arraybackslash}p{#1}}
\newcolumntype{Y}{>{\raggedright\arraybackslash}X}

% ---------- Datos del equipo (editar aqui una sola vez) ----------
\newcommand{\sigla}{ContX}
\newcommand{\organizacion}{Taller JC\textbar{}cell (Quevedo, Los R\'ios)}
\newcommand{\sistema}{Sistema de gesti\'on de reparaci\'on de equipos de c\'omputo}
\newcommand{\sistemaDiagrama}{de gesti\'on de\\reparaciones} % texto plano para el diagrama de contexto
\newcommand{\repourl}{https://github.com/EleSteven777/ISR401-ContX}

\setlength{\headheight}{14pt}
\pagestyle{fancy}
\fancyhf{}
\lhead{\small UTEQ $\cdot$ ISR-401}
\rhead{\small TA-PFC-E1 $\cdot$ ISR401-\sigla}
\cfoot{\small \thepage}

\begin{document}

% =====================================================================
%  PORTADA
% =====================================================================
\begin{titlepage}
  \centering
  {\Large\bfseries Universidad T\'ecnica Estatal de Quevedo\par}
  \vspace{0.2cm}
  {\large Facultad de Ciencias de la Computaci\'on\par}
  {\large Carrera de Software\par}
  \vspace{1.5cm}
  {\large Ingenier\'ia de Requerimientos (ISR-401) $\cdot$ 4.\textsuperscript{o} nivel $\cdot$ Paralelo A\par}
  \vspace{1.5cm}
  {\color{uteqverde}\rule{\linewidth}{1pt}\par}
  \vspace{0.4cm}
  {\LARGE\bfseries TA-PFC-E1\par}
  \vspace{0.3cm}
  {\Large Proyecto Fin de Curso, entrega 1:\\
   visi\'on, alcance, partes interesadas, contexto y plan de elicitaci\'on\par}
  \vspace{0.4cm}
  {\color{uteqverde}\rule{\linewidth}{1pt}\par}
  \vspace{1cm}
  {\large Organizaci\'on: \organizacion\par}
  {\large Sistema: \sistema\par}
  \vspace{1.5cm}
  {\large\bfseries Integrantes (equipo ISR401-\sigla)\par}
  \vspace{0.3cm}
  \begin{tabular}{ll}
    \pendiente{Apellidos, Nombres} & \pendiente{rol} \\
    \pendiente{Apellidos, Nombres} & \pendiente{rol} \\
    \pendiente{Apellidos, Nombres} & \pendiente{rol} \\
    \pendiente{Apellidos, Nombres} & \pendiente{rol} \\
  \end{tabular}
  \vfill
  {\large Docente: Gleiston Cicer\'on Guerrero Ulloa, PhD\par}
  \vspace{0.3cm}
  {\large Repositorio: \url{\repourl}\par}
  \vspace{0.3cm}
  {\large Quevedo, noviembre de 2026\par}
\end{titlepage}

\tableofcontents
\listoffigures
\listoftables
\clearpage

% =====================================================================
%  CUERPO (12 a 20 paginas de contenido)
% =====================================================================
% =====================================================================
\section{Introducci\'on}\label{sec:introduccion}
% =====================================================================
% Que debe ir aqui (aprox. 1 pagina):
%  - Problema de la organizacion en 1 parrafo (con datos de lo observado).
%  - Proposito de este documento y su lugar en el PFC.
%  - Marco de referencia: por que se identifican partes interesadas y el
%    limite del sistema ANTES de definir requisitos \cite{iso29148}; como
%    se eligen tecnicas de elicitacion segun la fuente y el tipo de
%    conocimiento \cite{pohl2025,ireb2026}.
%  - Organizacion del documento (una frase por seccion, con \ref).
% Redacte en tercera persona, sin relleno.

\pendiente{Redactar la introducci\'on con base en el problema real de la organizaci\'on.}
\pendiente{Marco: partes interesadas y l\'imite del sistema \cite{iso29148}; selecci\'on de t\'ecnicas \cite{pohl2025,ireb2026}.}

El documento se organiza as\'i: la Secci\'on~\ref{sec:organizacion} presenta la organizaci\'on;
la Secci\'on~\ref{sec:vision} define la visi\'on, los objetivos y el alcance;
la Secci\'on~\ref{sec:partes} identifica a las partes interesadas;
la Secci\'on~\ref{sec:contexto} describe el contexto del sistema;
la Secci\'on~\ref{sec:elicitacion} detalla el plan de elicitaci\'on;
la Secci\'on~\ref{sec:etica} trata la \'etica y el manejo de datos;
la Secci\'on~\ref{sec:gestion} presenta la gesti\'on del proyecto;
y la Secci\'on~\ref{sec:conclusiones} resume las conclusiones.

% =====================================================================
\section{Organizaci\'on}\label{sec:organizacion}
% =====================================================================

\subsection{Descripci\'on general}
El taller JC\textbar{}cell, ubicado en el cant\'on Quevedo, provincia de Los R\'ios, ofrece
mantenimiento y reparaci\'on de computadoras, laptops y otros dispositivos electr\'onicos.
Entre los servicios que mencion\'o el responsable est\'an el formateo, la instalaci\'on de
sistemas operativos y programas, la limpieza interna, el cambio de componentes y la soluci\'on
de fallas de software y hardware (SE-01, A2; EV-04; Anexo~\ref{anx:registro-entrevista}).

En la entrevista, el responsable estim\'o que el taller lleva cerca de tres a\~nos funcionando
(A2). Tambi\'en explic\'o que trabaja solo: recibe a los clientes, registra los equipos, realiza
los diagn\'osticos y reparaciones, comunica precios y coordina los cobros y las entregas (A1--A4,
B2--B8). Para registrar la recepci\'on anota el tipo, la marca, el modelo, la falla y los datos
de contacto; casi siempre lo hace en una libreta y a veces completa la informaci\'on en su
celular (B2). La antig\"uedad es una estimaci\'on del participante, no una fecha verificada.

\pendiente{Volumen aproximado de equipos atendidos por semana; no se pregunt\'o en SE-01.}

\subsection{Problema identificado}
En SE-01, el responsable describi\'o que consulta sus anotaciones para saber qu\'e equipos est\'an
en revisi\'on, pendientes de autorizaci\'on, en reparaci\'on o terminados; no cuenta con un sistema
que actualice esos estados (B2, B4; EV-04; Anexo~\ref{anx:registro-entrevista}). Cuando se
acumulan varios equipos, buscar entre las notas puede dificultar la identificaci\'on del equipo y
su estado. La principal dificultad que el participante prioriz\'o fue llevar el control de los
equipos que ingresan (C2, C4).

Como caso concreto, relat\'o que un equipo qued\'o esperando un repuesto entre varios trabajos
pendientes. D\'ias despu\'es tuvo que revisar las notas para recordar su estado y qu\'e faltaba hacer;
el tiempo de b\'usqueda le quit\'o tiempo a otras reparaciones. Indic\'o que esto puede ocurrir
cuando se acumulan equipos o una reparaci\'on tarda m\'as de lo previsto, pero no dio una frecuencia
ni un registro con el que medirla (C3). Tambi\'en mencion\'o dificultades para localizar equipos o
piezas, conseguir repuestos, avisar a clientes y recordar cobros pendientes (B2, B5--B7, C2).
Sobre los equipos que a veces no logra ubicar, la entrevista no aclara si se trata de un extrav\'io
temporal o una p\'erdida definitiva. Tampoco se revisaron registros que permitan confirmar esos
casos o establecer cu\'antos ocurren.

Lo anterior corresponde a lo declarado por el responsable en la entrevista. La transcripci\'on se
prepar\'o a partir de las notas del equipo y fue cotejada con la grabaci\'on autorizada (EV-04,
EV-05); aun as\'i, el audio registra el testimonio y no reemplaza los apuntes de recepci\'on o cobro
que permitir\'ian confirmar los casos y medir su frecuencia. El participante propuso revisar
ejemplos anonimizados de recepciones, trabajos y entregas para aclarar la informaci\'on pendiente
(E3). Esos ejemplos a\'un no constan entre las evidencias publicadas, por lo que las cifras de
frecuencia, tiempos de b\'usqueda y pagos olvidados siguen pendientes de contraste (C3, D1, E3).

\subsection{Primer contacto y evidencias}
La Tabla~\ref{tab:evidencias-org} resume las evidencias que respaldan la participaci\'on de la
organizaci\'on. Los originales se custodian en la carpeta institucional privada del equipo y su
registro (nombre de archivo, fecha y huella SHA-256) consta en \texttt{evidencias/MANIFIESTO.md}.

\begin{table}[H]
  \centering
  \caption{Evidencias de la relaci\'on con la organizaci\'on.}
  \label{tab:evidencias-org}
  \noindent\begin{tabularx}{\linewidth}{@{}L{1.2cm}L{3.5cm}L{2.2cm}Y@{}}
    \toprule
    \textbf{ID} & \textbf{Evidencia} & \textbf{Fecha} & \textbf{Descripci\'on} \\
    \midrule
    EV-02 & Registro del primer contacto (conversaci\'on previa)  & \pendiente{fecha} & \pendiente{medio y si existe registro escrito} \\
    EV-03 & Consentimiento informado firmado, PFC-ISR401-CONTX-A03 v1.3 (JC-731) & 20/11/2026 & Responsable del taller; modelo en blanco en el Anexo~\ref{anx:consentimiento} \\
    EV-04 & Transcripci\'on de la entrevista SE-01 (Word) & 21/11/2026 & Notas de la sesi\'on; registro anonimizado en el Anexo~\ref{anx:registro-entrevista} \\
    EV-05 & Grabaci\'on de audio de SE-01 & 21/11/2026 & 12 min 44 s; autorizada en la opci\'on b del consentimiento \\
    \bottomrule
  \end{tabularx}
\end{table}

% =====================================================================
\section{Visi\'on y alcance}\label{sec:vision}
% =====================================================================

\subsection{Declaraci\'on de visi\'on}
Para la persona responsable de JC\textbar{}cell, que necesita encontrar con facilidad los datos
de los equipos recibidos y conocer el avance de cada reparaci\'on, se propone una herramienta
sencilla que concentre el registro de ingreso, el diagn\'ostico y el estado del servicio. Se busca
reducir el tiempo que hoy dedica a revisar anotaciones y ayudarle a identificar los trabajos
pendientes, manteniendo su atenci\'on directa, sus diagn\'osticos y sus reparaciones. Esta es una
propuesta acad\'emica basada en SE-01 (B2, B4, C3--C4, D3; EV-04), no un compromiso de construir o
instalar el sistema; su alcance debe validarse con el responsable del taller.

\subsection{Objetivos de negocio}
El participante expres\'o que le gustar\'ia encontrar antes la informaci\'on de cada equipo, ubicar
las reparaciones pendientes y evitar confusiones y cobros olvidados (SE-01, C4, D1; EV-04).
La entrevista no incluy\'o mediciones iniciales ni metas acordadas. Por eso, los indicadores de la
Tabla~\ref{tab:objetivos} son propuestas para medir esos resultados; las l\'ineas base, metas y
plazos se completar\'an tras revisar ejemplos anonimizados y acordarlos con el responsable. Hasta
entonces, los objetivos siguen como borrador y no se consideran verificables.

\begin{table}[H]
  \centering
  \caption{Objetivos de negocio e indicadores propuestos para validar.}
  \label{tab:objetivos}
  \small
  \renewcommand{\arraystretch}{1.25}
  \noindent\begin{tabularx}{\linewidth}{@{}L{1.1cm}YL{2.6cm}L{2.4cm}L{2.4cm}@{}}
    \toprule
    \textbf{ID} & \textbf{Resultado buscado} & \textbf{Indicador} & \textbf{L\'inea base} & \textbf{Meta y plazo} \\
    \midrule
    ON-01 & Reducir el tiempo que toma localizar los datos y el estado de un servicio. Fuente: C3--C4, D1. & Mediana de minutos de las búsquedas completadas; registrar también las fallidas. & \pendiente{medir con ejemplos anonimizados (E3)} & \pendiente{acordar meta y fecha tras medir} \\
    ON-02 & Mejorar la consistencia del seguimiento de las reparaciones. Fuente: B4, C2, C4. & Servicios con estado coincidente / servicios verificables de la muestra $\times 100$. & \pendiente{revisar ejemplos de trabajos (E3)} & \pendiente{acordar porcentaje y plazo} \\
    ON-03 & Disminuir los cobros pendientes que quedan sin seguimiento al entregar un equipo. Fuente: B7, C2, D1. & Número de entregas con saldo pendiente confirmado, sin marca de pendiente en el registro. & \pendiente{revisar ejemplos de cobro (E3)} & \pendiente{acordar meta y fecha tras medir} \\
    \bottomrule
  \end{tabularx}
  \normalsize
\end{table}

Las fuentes de cada objetivo corresponden a preguntas de SE-01, recogidas en EV-04 y en el
Anexo~\ref{anx:registro-entrevista}. Los indicadores son propuestas del equipo, no mediciones
solicitadas o aprobadas por el participante.

Para medirlos, se propone acordar con el responsable el período, el tamaño de la muestra y los
criterios de selección antes de recoger los datos:
\begin{itemize}[nosep]
  \item \textbf{ON-01:} medir desde que se solicita localizar un servicio hasta encontrar sus
        datos y estado. Acordar un tiempo máximo y registrar las búsquedas fallidas por separado;
        la mediana se interpretará junto con la cantidad de búsquedas completadas y fallidas.
  \item \textbf{ON-02:} comparar el estado anotado con el que pueda reconstruirse junto al
        responsable, usando registros de trabajo o entrega del mismo momento de referencia.
        Los casos sin respaldo suficiente se contarán aparte, sin asumir que coinciden.
  \begin{samepage}
  \item \textbf{ON-03:} contrastar las notas del servicio y la entrega con la aclaración del
        responsable para identificar saldos pendientes que no quedaron marcados. Registrar solo
        el estado del cobro, sin datos de clientes; si no puede verificarse, dejar el caso pendiente.
  \par\end{samepage}
\end{itemize}
Las comparaciones posteriores usarán los mismos criterios; en ON-03 se usarán muestras del mismo
tamaño y se indicará el período revisado. La medición y sus condiciones aún deben acordarse y
realizarse; por eso se mantienen pendientes las líneas base, metas y fechas.

\subsection{Alcance}
La propuesta se centra en que el responsable pueda registrar y consultar los servicios del taller.
Las capacidades de la Tabla~\ref{tab:alcance} se derivan de SE-01 (B2--B7, C4, D1--D3, E2;
EV-04) y siguen pendientes de revisi\'on con el responsable; la entrevista no defini\'o todos los
detalles de operaci\'on.

\begin{table}[H]
  \centering
  \caption{Alcance propuesto para validar con el responsable.}
  \label{tab:alcance}
  \small
  \renewcommand{\arraystretch}{1.2}
  \noindent\begin{tabularx}{\linewidth}{@{}L{1.3cm}YL{1.6cm}L{2cm}@{}}
    \toprule
    \textbf{ID} & \textbf{Capacidad propuesta} & \textbf{Objetivo} & \textbf{Fuente} \\
    \midrule
    \multicolumn{4}{@{}l}{\textit{Dentro del alcance propuesto}} \\
    AL-01 & Registrar la recepci\'on del equipo, sus datos b\'asicos, la falla reportada y el contacto asociado. & ON-02 & B2, C4 \\
    AL-02 & Guardar el diagn\'ostico, la autorizaci\'on del cliente y el estado de cada reparaci\'on. & ON-02 & B3--B4 \\
    AL-03 & Consultar y localizar equipos y servicios activos, incluidos los que esperan revisi\'on, autorizaci\'on, repuesto o entrega. & ON-01 & B4--B5, C3--C4 \\
    AL-04 & Anotar el valor comunicado y si el servicio qued\'o pagado o pendiente, sin procesar pagos. & ON-03 & B3, B7 \\
    AL-05 & Consultar el historial de reparaciones asociadas a un equipo. & ON-01 & E2 \\
    \midrule
    \multicolumn{4}{@{}l}{\textit{Exclusiones propuestas, pendientes de acuerdo}} \\
    FA-01 & Envío automático de mensajes o integración con WhatsApp; se describieron llamadas y mensajes, pero no se definió su automatización. & --- & B6 \\
    FA-02 & Control de inventario, reposición automática o pedidos a proveedores; se mencionaron faltantes y pérdidas, pero no se definieron las funciones de ese módulo. & --- & B5, C2 \\
    FA-03 & Procesamiento de pagos, facturación electrónica o contabilidad; se describió el registro del cobro, sin definir esas funciones. & --- & B7 \\
    FA-04 & Administración de reglas y plazos de garantía; se mencionó su seguimiento, pero no se especificaron las condiciones. & --- & B8, E2 \\
    \bottomrule
  \end{tabularx}
  \normalsize
\end{table}

La columna de fuente remite a preguntas de SE-01 (EV-04;
Anexo~\ref{anx:registro-entrevista}). Los vínculos con los objetivos son propuestas del equipo;
por ejemplo, B7 respalda el problema de cobros de ON-03 y la capacidad AL-04. AL-05 se relaciona
con ON-01 porque permitiría consultar antecedentes de un equipo; ese vínculo aún debe validarse.

Las exclusiones son propuestas del equipo, pendientes de acuerdo con el taller. La falta de
detalle en la entrevista no demuestra que una función sea innecesaria. Antes de cerrar el alcance
se confirmarán tanto las inclusiones como las exclusiones con el responsable.

\subsection{Restricciones conocidas}
La Tabla~\ref{tab:restricciones} separa las condiciones expresadas por el responsable de los datos
que todav\'ia deben confirmarse. No se fija una plataforma t\'ecnica ni una obligaci\'on legal que no
haya sido verificada para este proyecto.

\begin{table}[H]
  \centering
  \caption{Restricciones y asuntos por confirmar.}
  \label{tab:restricciones}
  \small
  \renewcommand{\arraystretch}{1.2}
  \noindent\begin{tabularx}{\linewidth}{@{}L{1.1cm}L{2.2cm}YL{2.4cm}@{}}
    \toprule
    \textbf{ID} & \textbf{Tipo} & \textbf{Condici\'on} & \textbf{Origen} \\
    \midrule
    RE-01 & Datos y privacidad & Los ejemplos usados en el proyecto deben mantenerse anonimizados; los originales y datos identificables no se publican. El tratamiento, acceso y conservaci\'on de datos reales del sistema a\'un deben definirse. & Manifiesto de evidencias; EV-03 y EV-04 \\
    RE-02 & Recursos & La propuesta debe considerar la computadora, los celulares y la conexi\'on disponibles, y no requerir equipos costosos ni conocimientos t\'ecnicos avanzados. La plataforma y la calidad m\'inima de conexi\'on quedan por validar. & SE-01, D2--D3; EV-04 \\
    RE-03 & Plazo del proyecto & La fecha l\'imite y el calendario de hitos deben confirmarse con el cronograma oficial antes de fijar metas. & Cronograma del PFC; fecha pendiente de confirmaci\'on \\
    \bottomrule
  \end{tabularx}
  \normalsize
\end{table}

\section{Partes interesadas}\label{sec:partes}

Se consideran los roles que usan, aportan información o influyen en la propuesta, siguiendo la
identificación de fuentes de requisitos~\cite{ireb2026}. Los roles del taller se obtuvieron de
SE-01 (A1--A4, B2--B8; EV-04); los del equipo y el docente corresponden al proyecto académico.

\subsection{Registro de partes interesadas}
La Tabla~\ref{tab:partes} distingue ocho roles, no ocho personas. El responsable del taller
también realiza recepción, reparación y cobro (PI-01 a PI-04). No se supone que existan empleados
adicionales. Clientes y proveedores fueron mencionados, pero no entrevistados. Las valoraciones
de interés e influencia son una estimación del equipo.

\begingroup\small\setstretch{1.1}
\begin{longtable}{@{}L{1cm}L{2.6cm}L{4.5cm}L{1.3cm}L{1.5cm}L{2.4cm}@{}}
\caption{Roles relacionados con el proyecto.}\label{tab:partes}\\
\toprule
\textbf{ID} & \textbf{Rol} & \textbf{Interés y fundamento} & \textbf{Interés} & \textbf{Influen\-cia} & \textbf{Disponi\-bilidad} \\
\midrule\endfirsthead
\toprule
\textbf{ID} & \textbf{Rol} & \textbf{Interés y fundamento} & \textbf{Interés} & \textbf{Influen\-cia} & \textbf{Disponi\-bilidad} \\
\midrule\endhead
\bottomrule\endfoot
PI-01 & Responsable del taller & Decide prioridades y usaría la herramienta (A1, C4, D4). & Alto & Alta & Por acordar con el taller. \\
PI-02 & Recepción y entrega & Necesita identificar el equipo y registrar su ingreso y retiro (B2, B8). & Alto & Alta & La misma de PI-01. \\
PI-03 & Técnico de reparación & Necesita consultar diagnóstico, autorización y avance (B3--B5). & Alto & Alta & La misma de PI-01. \\
PI-04 & Registro y cobro & Necesita distinguir servicios pagados y pendientes (B7). & Alto & Alta & La misma de PI-01. \\
PI-05 & Cliente & Aporta la falla, autoriza trabajos y recibe el equipo (B1, B3, B6, B8). & Alto & Baja & Sin contacto directo; por coordinar. \\
PI-06 & Proveedor de repuestos & Informa disponibilidad y precio de piezas (A4, B5). & Medio & Baja & Sin contacto directo; por coordinar. \\
PI-07 & Equipo de requisitos & Analiza, documenta y revisa la propuesta. & Alto & Alta & Según acuerdos internos. \\
PI-08 & Docente evaluador & Define y evalúa los criterios académicos de la entrega. & Alto & Alta & Según horario académico; por confirmar. \\
\end{longtable}
\endgroup

\subsection{Matriz poder--interés}
La Figura~\ref{fig:matriz} representa la valoración inicial de la tabla. Los roles PI-01 a
PI-04 corresponden a una sola persona y no representan cuatro decisiones independientes.
La influencia se valora sobre esta propuesta académica, no sobre toda la actividad del taller.

\begin{figure}[H]
\centering
% Posiciones cualitativas propuestas; PI-01 a PI-04 son roles de la misma persona.
\begin{tikzpicture}[x=0.76cm,y=0.68cm,
parte/.style={rectangle,rounded corners=2pt,fill=uteqverde,text=white,inner sep=2pt,font=\scriptsize\bfseries}]
\fill[gray!6] (0,0) rectangle (10,10);
\fill[green!9] (5,5) rectangle (10,10);
\draw[thick] (0,0) rectangle (10,10);
\draw[dashed] (5,0)--(5,10);\draw[dashed] (0,5)--(10,5);
\node[font=\small\itshape,align=center] at (2.5,9.4) {Mantener satisfecho};
\node[font=\small\itshape,align=center] at (7.5,9.4) {Gestionar de cerca};
\node[font=\small\itshape] at (2.5,0.7) {Monitorear};
\node[font=\small\itshape] at (7.5,0.7) {Mantener informado};
\draw[-{Stealth}] (0,0)--(10.5,0) node[right] {Interés};
\draw[-{Stealth}] (0,0)--(0,10.5) node[above] {Poder};
\node[below,font=\scriptsize] at (0,0) {Bajo};
\node[below,font=\scriptsize] at (10,0) {Alto};
\node[left,font=\scriptsize] at (0,10) {Alto};
\node[parte] at (9.1,8.3) {PI-01};
\node[parte] at (6.1,7.1) {PI-02};
\node[parte] at (7.6,7.1) {PI-03};
\node[parte] at (9.1,7.1) {PI-04};
\node[parte] at (8.5,2.4) {PI-05};
\node[parte] at (5.8,2.4) {PI-06};
\node[parte] at (6.1,8.3) {PI-07};
\node[parte] at (7.6,8.3) {PI-08};
\end{tikzpicture}

\caption{Matriz poder--interés propuesta por el equipo.}\label{fig:matriz}
\end{figure}

El taller concentra las decisiones operativas. El equipo y el docente influyen en el contenido
y la calidad del trabajo académico. Clientes y proveedores aportan necesidades indirectas;
sus opiniones aún deben contrastarse si afectan decisiones del alcance.

\subsection{Estrategia de participación}
La Tabla~\ref{tab:estrategia} propone cómo involucrar a cada rol. Las consultas de PI-01 a
PI-04 se agruparán en una sola conversación para respetar el tiempo del responsable.

\begin{table}[H]\centering\small
\caption{Participación prevista por rol.}\label{tab:estrategia}
\noindent\begin{tabularx}{\linewidth}{@{}L{1cm}L{3cm}Y@{}}
\toprule
\textbf{ID} & \textbf{Estrategia} & \textbf{Acción y responsable} \\
\midrule
PI-01 & Gestionar de cerca & Jordan coordinará la revisión de prioridades y alcance antes de cerrarlos. \\
PI-02 & Gestionar de cerca & El equipo revisará un ejemplo de recepción y entrega en SE-02. \\
PI-03 & Gestionar de cerca & El equipo comprobará cómo se registra y consulta el estado en SE-02. \\
PI-04 & Gestionar de cerca & Jonathan revisará ejemplos sin datos personales para distinguir pagado y pendiente. \\
PI-05 & Mantener informado & Jordan propondrá el cuestionario CUE-01; participación y fecha por acordar. \\
PI-06 & Mantener informado & Se consultará disponibilidad de piezas solo si cambia el alcance; en E1 se usa la referencia indirecta de B5. \\
PI-07 & Gestionar de cerca & Los integrantes revisarán cambios y pendientes antes de cada versión. \\
PI-08 & Gestionar de cerca & El equipo contrastará la entrega con la guía y llevará dudas a tutoría. \\
\bottomrule
\end{tabularx}
\end{table}

\section{Contexto del sistema}\label{sec:contexto}

\subsection{Límite del sistema}
El límite separa la herramienta propuesta de las actividades que la rodean~\cite{ireb2026}.
Dentro quedan el registro y la consulta de equipos, diagnósticos, autorizaciones, estados,
cobros e historial (AL-01 a AL-05). Fuera quedan la reparación física, la compra de piezas,
las llamadas, los mensajes y el procesamiento de pagos. Esas actividades siguen a cargo del
responsable. Los archivos personales guardados en los equipos no son parte del estudio.

\subsection{Diagrama de contexto}
En la Figura~\ref{fig:contexto}, el responsable ingresa los datos y consulta los resultados.
Clientes y proveedores participan en el negocio, pero no se ha propuesto darles acceso directo
ni integrar WhatsApp o servicios externos. El diagrama representa el alcance propuesto, aún
pendiente de confirmar con el taller.

\begin{figure}[H]\centering
\resizebox{\linewidth}{!}{% Contexto de la propuesta: solo el responsable interactua directamente con el sistema.
\begin{tikzpicture}[
entidad/.style={rectangle,draw,thick,fill=gray!7,align=center,minimum width=3cm,minimum height=1.4cm},
sistema/.style={circle,draw=uteqverde,very thick,fill=green!8,align=center,minimum size=3.4cm},
flujo/.style={-{Stealth},thick},etq/.style={font=\small,align=center,fill=white,inner sep=2pt}]
\node[entidad] (R) at (0,0) {Responsable\\del taller\\(PI-01 a PI-04)};
\node[sistema] (S) at (7,0) {Sistema de gestión\\de reparaciones};
\draw[flujo] (R.north east) to[bend left=24] node[etq,above] {F1: datos del servicio\\y actualizaciones} (S.north west);
\draw[flujo] (R.east)--node[etq,above] {F2: consultas} (S.west);
\draw[flujo] (S.south west) to[bend left=24] node[etq,below] {F3: ficha, estado,\\pendientes e historial} (R.south east);
\end{tikzpicture}
}
\caption{Contexto de la herramienta propuesta.}\label{fig:contexto}
\end{figure}

La Tabla~\ref{tab:flujos} relaciona los flujos con las capacidades y las respuestas de SE-01.
\begin{table}[H]\centering\small
\caption{Flujos de información del contexto.}\label{tab:flujos}
\begin{tabularx}{\linewidth}{@{}L{1cm}YL{2.4cm}L{2.5cm}@{}}
\toprule
\textbf{Flujo} & \textbf{Contenido} & \textbf{Alcance} & \textbf{Fuente (EV-04)} \\
\midrule
F1 & Equipo, falla, contacto, diagnóstico, autorización, estado y condición del cobro. & AL-01, AL-02, AL-04 & B2--B4, B7 \\
F2 & Búsqueda de un equipo, servicio o antecedente. & AL-03, AL-05 & C3--C4, E2 \\
F3 & Datos encontrados, avance, trabajos pendientes e historial registrado. & AL-03, AL-05 & B4, C4, E2 \\
\bottomrule
\end{tabularx}
\end{table}

\section{Plan de elicitación}\label{sec:elicitacion}

\subsection{Criterios de selección de técnicas}
La entrevista recoge explicaciones que el responsable puede expresar. La observación ayudará
a reconocer pasos que realiza de forma habitual y quizá no menciona. El cuestionario permitirá
recoger la experiencia de clientes con preguntas breves. Se combinan fuentes y técnicas según
el conocimiento disponible~\cite{pohl2025,ireb2026}.

La Tabla~\ref{tab:tecnicas} distingue las fuentes consultadas de las actividades propuestas.
Solo se ha realizado la entrevista SE-01; las demás actividades requieren coordinación y el
consentimiento correspondiente.

\begingroup\small\setstretch{1.1}
\begin{longtable}{@{}L{1cm}L{2.7cm}L{2.4cm}L{5.7cm}L{1.8cm}@{}}
\caption{Técnicas y fuentes por rol.}\label{tab:tecnicas}\\
\toprule
\textbf{Rol} & \textbf{Técnica} & \textbf{Conoci\-miento} & \textbf{Motivo y estado} & \textbf{Instru\-mento} \\
\midrule\endfirsthead
\toprule
\textbf{Rol} & \textbf{Técnica} & \textbf{Conoci\-miento} & \textbf{Motivo y estado} & \textbf{Instru\-mento} \\
\midrule\endhead
\bottomrule\endfoot
PI-01 & Entrevista & Explícito: tareas y prioridades. & SE-01 realizada; permite explicar decisiones y casos concretos. & ENT-01 \\
PI-02 & Observación & Habitual o tácito: pasos de recepción y entrega. & Prevista para contrastar qué se anota y cómo se localiza el servicio. & OBS-01 \\
PI-03 & Observación & Habitual: actualización del estado. & Prevista; se observará el registro, sin intervenir equipos ni acceder a archivos privados. & OBS-01 \\
PI-04 & Revisión de ejemplos & Documentado: estado del cobro. & Prevista con registros anonimizados; solo se distinguirá pagado o pendiente. & OBS-01, O5 \\
PI-05 & Cuestionario & Explícito: experiencia del servicio. & Previsto; contrasta claridad de recepción, estado y entrega. No mide la línea base interna. & CUE-01 \\
PI-06 & Fuente indirecta & Declarado por el taller. & Se usa A4 y B5 de SE-01; no se consulta directamente mientras el módulo de inventario siga excluido. & ENT-01 \\
PI-07 & Revisión documental & Analítico: relaciones y decisiones. & El equipo contrasta entrevista, objetivos, modelos e instrumentos antes de integrar cambios. & Informe y EV-04 \\
PI-08 & Revisión de la guía & Normativo académico. & La guía define los criterios; las dudas se llevarán a tutoría, sin suponer aprobación previa. & Guía de E1 \\
\end{longtable}
\endgroup

\subsection{Cronograma de sesiones}
La Tabla~\ref{tab:sesiones} propone las siguientes actividades. Las semanas futuras son una
planificación del equipo; sus fechas dependen de la disponibilidad de quienes participen.

\begin{table}[H]\centering\small
\caption{Sesiones realizadas y previstas.}\label{tab:sesiones}
\noindent\begin{tabularx}{\linewidth}{@{}L{1cm}L{2.9cm}L{2.7cm}YL{2.5cm}@{}}
\toprule
\textbf{ID} & \textbf{Estado y momento} & \textbf{Rol / técnica} & \textbf{Propósito} & \textbf{Responsable} \\
\midrule
SE-01 & Realizada: 21/11/2026, 14:20; 12 min 44 s. & PI-01 / entrevista ENT-01. & Conocer proceso, dificultades y prioridades. & Jordan; notas de Jonathan. \\
SE-02 & Prevista: semana 5; fecha por acordar. & PI-02 a PI-04 / observación y ejemplos. & Contrastar registros, estados y datos para ON-01 a ON-03. & Jordan coordina; Jonathan registra. \\
SE-03 & Prevista: semana 5; fecha y participantes por acordar. & PI-05 / cuestionario. & Recoger experiencia del cliente sobre información y entrega. & Jordan coordina; equipo analiza. \\
SE-04 & Prevista: semana 6; fecha por acordar. & PI-01 / revisión conjunta. & Confirmar objetivos, alcance, exclusiones y correcciones. & Jonathan prepara; Jordan coordina. \\
\bottomrule
\end{tabularx}
\end{table}

SE-01 está resumida en el Anexo~\ref{anx:registro-entrevista}. Las nuevas sesiones requieren
explicar la actividad y sus permisos antes de comenzar. No se registran como resultados ya obtenidos.

\subsection{Instrumentos y trazabilidad}
Los anexos~\ref{anx:entrevista}, \ref{anx:cuestionario} y~\ref{anx:observacion} contienen los tres
instrumentos. La Tabla~\ref{tab:traza-instrumentos} muestra su relación con la propuesta actual.
La relación del guion con los objetivos se estableció después de analizar la entrevista; no
supone que esos objetivos ya estuvieran acordados al aplicarla.

\begin{table}[H]\centering\small
\caption{Relación de preguntas y observaciones con objetivos y alcance.}\label{tab:traza-instrumentos}
\noindent\begin{tabularx}{\linewidth}{@{}L{3.7cm}L{2.6cm}YL{2.4cm}@{}}
\toprule
\textbf{Preguntas / aspectos} & \textbf{Objetivo} & \textbf{Capacidad} & \textbf{Estado} \\
\midrule
ENT-01: B2, B4, C3--C4, D1 & ON-01, ON-02 & AL-01 a AL-03 & Aplicado. \\
ENT-01: B3, B7, C2, D1 & ON-03 & AL-04 & Aplicado. \\
ENT-01: E2 & ON-01 & AL-05 & Aplicado; alcance por confirmar. \\
CUE-01: C1--C4 & ON-02 & AL-01 a AL-03; percepción del cliente. & Sin aplicar. \\
CUE-01: C5 & ON-03 & AL-04; claridad del estado de pago. & Sin aplicar. \\
OBS-01: O1--O4 & ON-01, ON-02 & AL-01 a AL-03 & Sin aplicar. \\
OBS-01: O5 & ON-03 & AL-04 & Sin aplicar. \\
OBS-01: O6 & ON-01 & AL-05 & Sin aplicar. \\
\bottomrule
\end{tabularx}
\end{table}

% =====================================================================
\section{\'Etica y manejo de datos}\label{sec:etica}
% =====================================================================
% Fundamento: LOPDP \cite{lopdp2021}, su Reglamento \cite{rlopdp2023} y
% el Codigo de Etica de la ACM \cite{acm2018}. Cite articulos solo si
% los ha leido en el texto oficial y dicen lo que usted afirma.

El tratamiento de los datos de los participantes se rige por la Ley Org\'anica de Protecci\'on de
Datos Personales \cite{lopdp2021} y su Reglamento \cite{rlopdp2023}; la conducta del equipo, por
el C\'odigo de \'Etica de la ACM \cite{acm2018}.

\subsection{Consentimiento informado}
Antes de cada sesi\'on, cada participante recibe y firma el modelo de consentimiento del
Anexo~\ref{anx:consentimiento}. Los consentimientos firmados no se publican en el repositorio.

\pendiente{Describir c\'omo se entrega, explica y archiva el consentimiento; qu\'e ocurre si un participante lo revoca.}

\subsection{Anonimizaci\'on}
% Ejemplo de procedimiento: codigos P01, P02... por participante; la
% tabla codigo-identidad vive solo en la carpeta privada; transcripciones
% sin nombres, cedulas, telefonos ni correos; voces/rostros no se suben.
\pendiente{Procedimiento paso a paso de anonimizaci\'on aplicado por el equipo.}

\subsection{Custodia de evidencias}
Las evidencias originales (carta de aceptaci\'on, consentimientos firmados, audios, fotograf\'ias)
se guardan en una carpeta institucional privada (cuenta \texttt{@uteq.edu.ec}) compartida con el
docente. En el repositorio solo se registra el archivo \texttt{evidencias/MANIFIESTO.md} con el
nombre de archivo, la fecha, la duraci\'on (si aplica) y la huella SHA-256 de cada evidencia, lo
que permite comprobar que el original no fue alterado.

\pendiente{Indicar qui\'en del equipo administra la carpeta, plazo de conservaci\'on y eliminaci\'on.}

\section{Gestión del proyecto}\label{sec:gestion}


\subsection{Aplicación de los principios del PMBOK}


La organización del equipo se apoya en los principios del PMBOK, séptima edición~\cite{pmbok2021}. Se asignan responsabilidades principales y se mantiene la revisión conjunta de los entregables. Esta distribución adapta el trabajo al tamaño del equipo y al contexto académico del proyecto; sus denominaciones no son cargos obligatorios del PMBOK.

\begin{itemize}
\item \textbf{Equipo:} los tres integrantes colaboran y conocen el contenido de los entregables, aunque cada uno coordina actividades diferentes.
\item \textbf{Interesados:} Jordan coordina las sesiones y la comunicación con el taller. Las interpretaciones que requieran confirmación se consultarán con los participantes pertinentes.
\item \textbf{Calidad:} Rai prepara los documentos y modelos; Jonathan revisa su claridad, coherencia y respaldo antes de cerrar cada versión.
\item \textbf{Adaptación al contexto:} el equipo ajustará las actividades a la disponibilidad del taller y a la información que pueda obtenerse, registrando los cambios acordados.
\item \textbf{Administración responsable:} Jordan administra la carpeta privada; todos respetan los consentimientos y evitan publicar información identificable.
\end{itemize}

Estas acciones constituyen una aplicación prevista de los principios, que deberá comprobarse mediante los aportes, revisiones, acuerdos y evidencias del proyecto.


\subsection{Roles y responsabilidades}


\begin{table}[H]
\centering\small
\caption{Roles y responsabilidades del equipo ContX.}\label{tab:roles}
\noindent\begin{tabularx}{\linewidth}{@{}L{3.6cm}L{3.2cm}Y@{}}
\toprule
\textbf{Integrante} & \textbf{Responsabilidad principal} & \textbf{Actividades} \\
\midrule
Holguin Mendoza Jonathan Alexander & Validación y trazabilidad & Revisar documentos y modelos; comprobar claridad, consistencia y verificabilidad; relacionar hallazgos, objetivos, alcance y evidencias; registrar observaciones y verificar las correcciones. \\
Alvarez Rubio Rai Eduardo & Análisis y documentación & Organizar y analizar la información obtenida; redactar los documentos; elaborar los modelos e incorporar las correcciones acordadas. \\
Basurto Anchaguano Jordan Steven & Coordinación y elicitación & Coordinar el contacto y las sesiones con el taller; preparar y conducir la elicitación con apoyo del equipo; gestionar aceptación y consentimientos; custodiar las evidencias privadas. \\
\bottomrule
\end{tabularx}
\end{table}


Las responsabilidades no son exclusivas: todos contribuyen y revisan el trabajo.

\subsection{Procedimiento de revisión y validación}
\label{sec:revision-entregables}

El responsable de validación y trazabilidad coordinará la revisión antes de cada entrega o actualización de un instrumento, documento o modelo:

\begin{enumerate}
\item \textbf{Delimitar la revisión.} Identificar el archivo y su versión, las fuentes disponibles y los productos que solicita la guía.
\item \textbf{Revisar el contenido.} Comprobar que las preguntas sean claras, neutrales y coherentes con el consentimiento; que visión, objetivos, alcance, partes interesadas, técnicas e instrumentos concuerden; y que los modelos y las conclusiones tengan respaldo. En los objetivos, revisar indicador, situación inicial, meta y plazo.
\item \textbf{Comprobar la trazabilidad.} Verificar las relaciones entre fuentes, hallazgos, objetivos, alcance y modelos según el criterio descrito a continuación.
\item \textbf{Comunicar observaciones.} Indicar el elemento afectado, el problema y la corrección propuesta. Separar los ajustes documentales de las dudas que deben consultarse al taller y revisar que no se publiquen datos personales.
\item \textbf{Verificar las correcciones.} El responsable del contenido realiza los ajustes y el de validación comprueba su resolución. Si cambian las necesidades, prioridades o el alcance, Jordan coordina su confirmación con el taller; la revisión interna no sustituye ese acuerdo.
\item \textbf{Cerrar la revisión.} Registrar lo revisado, las correcciones y los pendientes. Comprobar las citas bibliográficas, la compilación y las referencias a tablas y figuras antes de integrar la versión al repositorio.
\end{enumerate}

Las observaciones se marcarán como pendientes, en corrección o resueltas. Los datos y acuerdos sin confirmar conservarán su estado pendiente.

\subsection{Criterio de trazabilidad para E1}

Se usarán los identificadores de preguntas, instrumentos, partes interesadas, objetivos, alcance y evidencias del documento. Cada relación indicará su fuente, la referencia concreta y el estado de confirmación, mediante las tablas, figuras y referencias del informe. El manifiesto permitirá identificar los originales privados.

La relación será: \textbf{fuente y evidencia $\rightarrow$ hallazgo $\rightarrow$ objetivo o capacidad del alcance $\rightarrow$ elemento del modelo}, cuando corresponda. Un hallazgo también puede justificar una restricción o exclusión. Los requisitos detallados se vincularán cuando se elaboren en las siguientes entregas.

\subsection{Cronograma hasta la semana 17}
% Semana 4 = 16-22 nov 2026. Marque con \x las semanas de cada actividad.
% Ajuste las actividades a las entregas reales del PFC del silabo.
La Tabla~\ref{tab:cronograma} propone la distribución del trabajo hasta la semana 17.
Las semanas de las actividades posteriores a E1 son una planificación del equipo y se ajustarán
a las fechas oficiales de las siguientes entregas.

\newcommand{\x}{\cellcolor{uteqverde!60}}
\begin{table}[H]
  \centering
  \caption{Cronograma del PFC (semanas 4 a 17).}
  \label{tab:cronograma}
  \small
  \setlength{\tabcolsep}{3pt}
  \begin{tabular}{@{}L{5.2cm}*{14}{|c}|@{}}
    \toprule
    \textbf{Actividad} & 4 & 5 & 6 & 7 & 8 & 9 & 10 & 11 & 12 & 13 & 14 & 15 & 16 & 17 \\
    \midrule
    E1: visión, alcance, contexto y plan & \x &  &  &  &  &  &  &  &  &  &  &  &  &  \\
    Completar fuentes y mediciones &  & \x & \x &  &  &  &  &  &  &  &  &  &  &  \\
    Analizar y especificar requisitos &  &  & \x & \x & \x & \x &  &  &  &  &  &  &  &  \\
    Modelar y revisar requisitos &  &  &  &  &  & \x & \x & \x & \x &  &  &  &  &  \\
    Validar con la organización &  &  & \x &  &  &  & \x &  &  &  & \x & \x &  &  \\
    Corregir y consolidar el informe &  &  &  &  &  &  &  &  &  & \x & \x & \x & \x &  \\
    Integrar la entrega final del PFC &  &  &  &  &  &  &  &  &  &  &  &  &  & \x \\
    \bottomrule
  \end{tabular}
\end{table}

La entrevista SE-01 corresponde a la semana 4; las sesiones de las semanas siguientes
son una previsión y aún deben acordarse con el taller.

\subsection{Plan de comunicaci\'on con la organizaci\'on}
La Tabla~\ref{tab:comunicacion} define qu\'e se comunica, a qui\'en, por qu\'e medio y con qu\'e frecuencia.

\begin{table}[H]
  \centering
  \caption{Plan de comunicaci\'on.}
  \label{tab:comunicacion}
  \noindent\begin{tabularx}{\linewidth}{@{}YL{2.4cm}L{2.4cm}L{2.2cm}L{2.4cm}@{}}
    \toprule
    \textbf{Qu\'e se comunica} & \textbf{Destinatario} & \textbf{Medio} & \textbf{Frecuencia} & \textbf{Responsable} \\
    \midrule
    Avance y dudas por resolver & Responsable del taller (PI-01) & Resumen breve; canal por acordar & Al cerrar una versión o surgir una decisión & Jordan coordina; equipo prepara \\
    Agenda y condiciones de las sesiones & Roles convocados; PI-01 como contacto del taller & Canal por acordar, sin datos personales & Antes de cada actividad & Jordan \\
    Validación del problema, prioridades, objetivos y alcance propuesto & Responsable del taller (PI-01) & Por acordar con el taller & Antes de cerrar la propuesta y ante cambios relevantes & Jordan y Jonathan \\
    \bottomrule
  \end{tabularx}
\end{table}

Para esta consulta, Jonathan preparará un resumen de la propuesta y los puntos por confirmar.
Jordan coordinará la fecha y el medio con el responsable del taller. Se registrarán los acuerdos,
las correcciones y las decisiones pendientes, indicando la fecha y la versión revisada. La
actividad queda prevista; los resultados se incorporarán después de realizarla.

\subsection{Registro de riesgos}\label{sec:registro-riesgos}

La Tabla~\ref{tab:riesgos} recoge los riesgos que pueden afectar la elaboración y validación de los requisitos. Se basa en la entrevista SE-01 y en los datos que aún faltan por confirmar. Las respuestas son acciones previstas.

Se propone valorar la probabilidad (P) y el impacto (I) de 1 a 3: bajo, medio y alto. El producto P$\times$I servirá para ordenar la atención. Las valoraciones son estimaciones iniciales para revisión del equipo.

\begingroup
\small
\setstretch{1.1}
\setlength{\tabcolsep}{4pt}
\renewcommand{\arraystretch}{1.2}
\begin{longtable}{@{}L{0.9cm}L{4.5cm}cccL{4.7cm}L{2.5cm}@{}}
\caption{Riesgos del proyecto y respuestas previstas.}\label{tab:riesgos}\\
\toprule
\textbf{ID} & \textbf{Riesgo} & \textbf{P} & \textbf{I} & \textbf{P$\times$I} & \textbf{Respuesta} & \textbf{Responsable} \\
\midrule
\endfirsthead
\multicolumn{7}{@{}l}{\textit{Tabla~\ref{tab:riesgos}. Continuación.}}\\
\toprule
\textbf{ID} & \textbf{Riesgo} & \textbf{P} & \textbf{I} & \textbf{P$\times$I} & \textbf{Respuesta} & \textbf{Responsable} \\
\midrule
\endhead
\bottomrule
\endfoot
RI-01 & Formular objetivos que no puedan comprobarse, porque aún no hay cifras sobre la frecuencia de las dificultades ni el tiempo de búsqueda de información (C3 y D1). & 3 & 3 & 9 & Revisar los registros de ejemplo propuestos en E3 y acordar cómo obtener los datos que faltan. No completar cifras por suposición. Si no se obtienen, mantener la medición pendiente y explicar la limitación. & Análisis y documentación; revisión de validación y trazabilidad. \\
\addlinespace[7pt]
RI-02 & Retrasar las aclaraciones o la validación por depender del responsable del taller, quien también realiza toda la atención y reparación (A3 y D4). & 2 & 2 & 4 & Agrupar las dudas y coordinar con anticipación una conversación breve. Si no puede atenderla, continuar con los datos confirmados y dejar las decisiones pendientes para otra fecha acordada. & Coordinación y elicitación. \\
\addlinespace[7pt]
RI-03 & Convertir las dificultades mencionadas en funciones sin acordarlas con el taller, ampliando el alcance más allá de su prioridad de registro y control de equipos (C2 y C4 de SE-01). & 2 & 3 & 6 & Relacionar cada capacidad con una respuesta o evidencia y acordar con el responsable qué se incluirá. No presentar como confirmado lo que solo sea una propuesta del equipo. & Validación y trazabilidad; apoyo de análisis y documentación. \\
\addlinespace[7pt]
RI-04 & Publicar por descuido datos personales presentes en las notas, el consentimiento o los registros de ejemplo del taller, afectando la privacidad de participantes y clientes. & 2 & 3 & 6 & Mantener los originales en la carpeta privada y revisar la copia anonimizada antes de subirla. Si se detecta una publicación indebida, retirar el acceso público al archivo, avisar al docente y coordinar su tratamiento. & Custodia de evidencias; revisión de validación y trazabilidad. \\
\end{longtable}
\endgroup

El equipo revisará los riesgos durante el proyecto. Si alguno ocurre, el responsable aplicará la respuesta prevista.

\section{Conclusiones}\label{sec:conclusiones}

La entrevista identifica como prioridad organizar el registro y seguimiento de los equipos.
El caso de una reparación detenida por un repuesto muestra la dificultad de reconstruir su
estado entre varias anotaciones. Esto respalda una propuesta centrada en localizar servicios
y consultar su avance, manteniendo la atención directa del taller.

El alcance distingue registro de reparaciones y cobros de funciones más amplias, como inventario,
mensajería automática o facturación. Estas inclusiones y exclusiones son propuestas del equipo
y deben confirmarse con el responsable.

Se cuenta con consentimiento, transcripción y audio registrados en el manifiesto. La entrevista
aporta testimonio, pero no una medición de frecuencia, tiempos o pagos omitidos. El siguiente
paso es revisar ejemplos anonimizados, acordar metas y validar la propuesta; el cuestionario y
la observación están preparados, todavía sin aplicar.

% =====================================================================
\section*{Declaraci\'on de uso de inteligencia artificial}\label{sec:ia}
\addcontentsline{toc}{section}{Declaraci\'on de uso de inteligencia artificial}
% =====================================================================
% Obligatoria (sin ella: -0,3). Declare TODA herramienta usada, incluida
% la que haya servido para armar la estructura del repositorio o la
% plantilla LaTeX. Si no se uso IA en una parte, dilo expresamente.

La Tabla~\ref{tab:ia} declara las herramientas de inteligencia artificial usadas como apoyo.
El apoyo incluyó propuestas de redacción y revisión de instrumentos, ética, gestión y riesgos,
además de ajustes en los archivos y la compilación. Por ello, estos textos no se presentan como
redactados íntegramente sin IA. La revisión documental no sustituye la comprobación de los
originales, la revisión del equipo ni la validación con la organización.

\begin{table}[H]
  \centering
  \caption{Uso de herramientas de inteligencia artificial.}
  \label{tab:ia}
  \noindent\begin{tabularx}{\linewidth}{@{}L{2.6cm}L{2cm}YY@{}}
    \toprule
    \textbf{Herramienta} & \textbf{Versi\'on} & \textbf{Prop\'osito} & \textbf{Verificaci\'on} \\
    \midrule
    Codex (OpenAI) & GPT-6 en esta revisión; \pendiente{versiones de sesiones anteriores} & Apoyo en redacción, revisión de coherencia y contraste con la guía; edición de fuentes y compilación. & Se contrastaron fechas y roles con la información del equipo y se revisaron fuentes y compilación. \pendiente{revisión del equipo sobre los textos asistidos y cotejo de evidencias privadas} \\
    \bottomrule
  \end{tabularx}
\end{table}


% =====================================================================
%  REFERENCIAS (estilo IEEE)
% =====================================================================
\clearpage
\bibliographystyle{IEEEtran}
\bibliography{referencias}

% =====================================================================
%  ANEXOS: instrumentos y consentimiento
% =====================================================================
\clearpage
\appendix
\section{Guion de entrevista}\label{anx:entrevista}
% =====================================================================
%  Instrumento 1: guion de entrevista semiestructurada
%  Se incluye como anexo desde main.tex.
%  Cada pregunta debe trazarse a un objetivo de negocio (ON-xx) o a una
%  capacidad del alcance (AL-xx). Escriba preguntas abiertas, propias de
%  la organizacion; evite preguntas genericas validas para cualquier
%  sistema.
% =====================================================================

\subsection*{Datos de la sesi\'on}
\begin{tabularx}{\linewidth}{@{}L{4.5cm}Y@{}}
  \toprule
  C\'odigo de sesi\'on & SE-01 \\
  Parte interesada (rol) & \pendiente{PI-xx, rol} \\
  C\'odigo del participante & P\_\_\_ (sin nombre) \\
  Fecha, hora y lugar & \\
  Entrevistador / anotador & \\
  Duraci\'on prevista & \pendiente{45 min} \\
  Consentimiento firmado & S\'i \quad No \\
  \bottomrule
\end{tabularx}

\subsection*{Objetivo de la entrevista}
\pendiente{Qu\'e conocimiento se busca obtener de este rol y para qu\'e objetivo de negocio.}

\subsection*{Apertura (5 min)}
\begin{enumerate}[label=A\arabic*.]
  \item Presentaci\'on del equipo y del prop\'osito de la sesi\'on.
  \item Lectura del consentimiento; confirmaci\'on de grabaci\'on (s\'olo si el participante acepta).
\end{enumerate}

\subsection*{Bloque 1: rol y tareas actuales}
\begin{tabularx}{\linewidth}{@{}L{1cm}YL{1.6cm}@{}}
  \toprule
  \textbf{N.\textsuperscript{o}} & \textbf{Pregunta} & \textbf{Traza} \\
  \midrule
  1.1 & \pendiente{pregunta} & \pendiente{ON/AL} \\
  1.2 & \pendiente{pregunta} & \pendiente{ON/AL} \\
  1.3 & \pendiente{pregunta} & \pendiente{ON/AL} \\
  \bottomrule
\end{tabularx}

\subsection*{Bloque 2: problemas y causas}
\begin{tabularx}{\linewidth}{@{}L{1cm}YL{1.6cm}@{}}
  \toprule
  \textbf{N.\textsuperscript{o}} & \textbf{Pregunta} & \textbf{Traza} \\
  \midrule
  2.1 & \pendiente{pregunta} & \pendiente{ON/AL} \\
  2.2 & \pendiente{pregunta} & \pendiente{ON/AL} \\
  2.3 & \pendiente{pregunta} & \pendiente{ON/AL} \\
  \bottomrule
\end{tabularx}

\subsection*{Bloque 3: expectativas y restricciones}
\begin{tabularx}{\linewidth}{@{}L{1cm}YL{1.6cm}@{}}
  \toprule
  \textbf{N.\textsuperscript{o}} & \textbf{Pregunta} & \textbf{Traza} \\
  \midrule
  3.1 & \pendiente{pregunta} & \pendiente{ON/AL} \\
  3.2 & \pendiente{pregunta} & \pendiente{ON/AL} \\
  3.3 & \pendiente{pregunta} & \pendiente{ON/AL} \\
  \bottomrule
\end{tabularx}

\subsection*{Cierre (5 min)}
\begin{enumerate}[label=C\arabic*.]
  \item ?`Hay algo importante sobre su trabajo que no se haya preguntado?
  \item ?`Qu\'e otra persona (rol) deber\'ia entrevistarse?
  \item Agradecimiento y explicaci\'on de los siguientes pasos.
\end{enumerate}


\clearpage
\section{Cuestionario}\label{anx:cuestionario}
\subsection*{Identificación e instrucciones}
\textbf{Código:} CUE-01, versión 1.0. \quad
\textbf{Dirigido a:} clientes (PI-05). \quad
\textbf{Duración estimada:} 5 minutos.

Instrumento preparado, todavía sin aplicar. Antes de responder se explicará la actividad y
se solicitará el consentimiento correspondiente. No escriba nombres, teléfonos, números de
equipo ni montos. Puede omitir preguntas o terminar cuando desee. Piense en su servicio más reciente.

\subsection*{Experiencia del servicio}
Marque una respuesta por fila: \textbf{Sí}, \textbf{En parte}, \textbf{No} o
\textbf{No aplica / no recuerda}. Las preguntas se refieren a lo que pudo conocer como cliente,
no a los registros internos del taller.

\medskip
\noindent\begin{tabularx}{\linewidth}{@{}L{0.9cm}Ycccc@{}}
\toprule
\textbf{ID} & \textbf{Pregunta} & \textbf{Sí} & \textbf{En parte} & \textbf{No} & \textbf{N/A} \\
\midrule
C1 & Al dejar el equipo, ¿pudo confirmar qué problema se registró? & $\square$ & $\square$ & $\square$ & $\square$ \\
C2 & Cuando consultó por el servicio, ¿recibió información sobre su estado? & $\square$ & $\square$ & $\square$ & $\square$ \\
C3 & Si se necesitó un trabajo adicional, ¿le pidieron autorización antes de hacerlo? & $\square$ & $\square$ & $\square$ & $\square$ \\
C4 & ¿Le avisaron cuando el equipo estuvo listo para retirar? & $\square$ & $\square$ & $\square$ & $\square$ \\
C5 & Al retirar el equipo, ¿quedó claro si el servicio estaba pagado o pendiente? & $\square$ & $\square$ & $\square$ & $\square$ \\
\bottomrule
\end{tabularx}

\subsection*{Comentario opcional}
\textbf{C6.} ¿Qué información le habría ayudado durante el servicio? No incluya datos personales.

\medskip\noindent\rule{\linewidth}{0.4pt}\par\medskip
\noindent\rule{\linewidth}{0.4pt}

\medskip
El equipo contará las respuestas por opción y separará las omitidas. No presentará este
cuestionario como una medición de tiempos de búsqueda o de pagos omitidos dentro del taller.
La trazabilidad de C1--C5 consta en la Tabla~\ref{tab:traza-instrumentos}; C6 ayuda a detectar
necesidades que deban consultarse antes de cambiar el alcance.


\clearpage
\section{Gu\'ia de observaci\'on}\label{anx:observacion}
% =====================================================================
%  Instrumento 3: guia de observacion (de campo)
%  Se incluye como anexo desde main.tex.
%  Sirve para conocimiento tacito/subconsciente: lo que la gente hace sin
%  poder explicarlo. Defina QUE proceso se observa y QUE aspectos se
%  registran; no fotografie personas ni documentos con datos personales.
% =====================================================================

\subsection*{Datos de la observaci\'on}
\noindent\begin{tabularx}{\linewidth}{@{}L{4.5cm}Y@{}}
  \toprule
  C\'odigo de sesi\'on & SE-02 \\
  Proceso observado & \pendiente{proceso} \\
  Rol(es) observado(s) & \pendiente{PI-xx} \\
  Fecha, hora inicio / fin & \\
  Lugar & \\
  Observador(es) & \\
  Tipo de observaci\'on & $\square$ No participante \quad $\square$ Participante \\
  Consentimiento & $\square$ Firmado por los observados \\
  \bottomrule
\end{tabularx}

\subsection*{Objetivo}
\pendiente{Qu\'e se quiere comprender del proceso y con qu\'e objetivo de negocio se relaciona.}

\subsection*{Aspectos a observar}
\noindent\begin{tabularx}{\linewidth}{@{}L{0.8cm}L{4.2cm}YL{1.5cm}@{}}
  \toprule
  \textbf{N.\textsuperscript{o}} & \textbf{Aspecto} & \textbf{Registro (qu\'e se vio)} & \textbf{Traza} \\
  \midrule
  O1 & Secuencia de pasos de la tarea & & \pendiente{AL} \\
  O2 & Documentos, formularios o herramientas usados & & \pendiente{AL} \\
  O3 & Tiempo por tarea / tiempos de espera & & \pendiente{ON} \\
  O4 & Errores, reprocesos o excepciones & & \pendiente{ON} \\
  O5 & Interrupciones y comunicaciones con otros roles & & \pendiente{AL} \\
  O6 & \pendiente{aspecto propio de la organizaci\'on} & & \pendiente{} \\
  \bottomrule
\end{tabularx}

\subsection*{Registro cronol\'ogico}
\noindent\begin{tabularx}{\linewidth}{@{}L{1.6cm}YL{4cm}@{}}
  \toprule
  \textbf{Hora} & \textbf{Hecho observado} & \textbf{Comentario del observador} \\
  \midrule
  & & \\[0.5cm]
  & & \\[0.5cm]
  & & \\[0.5cm]
  & & \\[0.5cm]
  \bottomrule
\end{tabularx}

\subsection*{Preguntas de aclaraci\'on posteriores}
\begin{enumerate}[label=Q\arabic*.]
  \item \rule{0.9\linewidth}{0.4pt}
  \item \rule{0.9\linewidth}{0.4pt}
\end{enumerate}


\clearpage
\section{Modelo de consentimiento informado}\label{anx:consentimiento}
% =====================================================================
%  Modelo de consentimiento informado (PFC-ISR401-REPCOMP-A01, v1.0)
%  Se incluye como anexo desde main.tex. El ejemplar FIRMADO nunca se
%  sube al repositorio: va a la carpeta institucional privada y se
%  registra en evidencias/MANIFIESTO.md con su huella SHA-256.
% =====================================================================

\begin{center}
  {\small\bfseries FACULTAD DE CIENCIAS DE LA COMPUTACI\'ON $\cdot$ CARRERA DE INGENIER\'IA DE SOFTWARE\par}
  \vspace{0.1cm}
  {\large\bfseries FORMULARIO DE CONSENTIMIENTO INFORMADO\par}
  {\small Participaci\'on voluntaria en el levantamiento de requisitos de un proyecto acad\'emico\par}
\end{center}

\begin{tabularx}{\linewidth}{@{}L{4cm}Y@{}}
  \toprule
  \textbf{C\'odigo del documento} & PFC-ISR401-REPCOMP-A01 \quad Versi\'on 1.0 (16/11/2026) \\
  \textbf{Asignatura} & Ingenier\'ia de Requisitos (ISR-401) $\cdot$ 4.\textsuperscript{o} nivel $\cdot$ Paralelo A $\cdot$ Per\'iodo 2026--2027 \\
  \textbf{Proyecto} & Sistema de gesti\'on de reparaci\'on de equipos de c\'omputo \\
  \textbf{Organizaci\'on estudiada} & Taller de reparaci\'on de equipos de c\'omputo (nombre: JC\textbar{}cell), cant\'on Quevedo, provincia de Los R\'ios \\
  \textbf{Docente responsable} & Ing. Gleiston Cicer\'on Guerrero Ulloa, PhD $\cdot$ \texttt{gguerrero@uteq.edu.ec} \\
  \textbf{Contacto del equipo} & Basurto Anchaguano Jordan Steven $\cdot$ \texttt{jbasurtoa@uteq.edu.ec} \\
  \textbf{C\'odigo del participante} & \textbf{REPCOMP-P01} \\
  \bottomrule
\end{tabularx}

\medskip
Usted ha sido invitado o invitada a participar en una actividad acad\'emica de la Universidad
T\'ecnica Estatal de Quevedo (UTEQ). Antes de decidir, lea con calma este documento o pida que se
lo lean en voz alta. Puede hacer todas las preguntas que necesite. Participar es opcional: si no
desea hacerlo, no tiene que dar ninguna explicaci\'on.

\subsection*{1. ?`Para qu\'e es este estudio?}
Un grupo de estudiantes de Ingenier\'ia de Software est\'a estudiando c\'omo se trabaja en un taller
de reparaci\'on de equipos de c\'omputo, en particular en la recepci\'on y registro de equipos, el
diagn\'ostico y la cotizaci\'on, el seguimiento de las reparaciones, el control de repuestos e
inventario, la comunicaci\'on con los clientes y la facturaci\'on y entrega, con el fin de
identificar, analizar y documentar los requisitos de un sistema inform\'atico (Sistema de gesti\'on
de reparaci\'on de equipos de c\'omputo). La informaci\'on se usar\'a como evidencia para elaborar
documentos de ingenier\'ia de requisitos: listas de requisitos, casos de uso, modelos, prototipos y
registros de validaci\'on. Se trata de una propuesta acad\'emica: su participaci\'on no implica la
compra, instalaci\'on ni entrega de ning\'un sistema a la organizaci\'on.

\subsection*{2. ?`Qu\'e tendr\'ia que hacer?}
Seg\'un su rol (responsable del taller), se le pedir\'a una o varias de estas actividades, que el
equipo le indicar\'a antes de empezar:
\begin{itemize}[nosep]
  \item Una entrevista sobre sus tareas diarias, los problemas que encuentra y las mejoras que
        considera necesarias (aproximadamente 15 a 45 minutos).
  \item Responder a un formulario breve ($\square$ impreso o $\square$ en l\'inea; marque el que prefiera).
  \item Permitir que un integrante del equipo observe su trabajo sin intervenir (observaci\'on no participante).
  \item Revisar las interfaces de los prototipos del posible sistema y opinar si reflejan su trabajo
        (sesi\'on de validaci\'on).
\end{itemize}
Lo que usted diga se registrar\'a con notas escritas. El audio, el video y las fotograf\'ias solo se
aplicar\'an si usted lo autoriza en la Secci\'on 7. Negarse a ser grabado no le impide participar;
sin embargo, estos registros son de gran importancia en nuestro trabajo.

\textbf{No se le pedir\'a:} usuarios ni contrase\~nas, ni contenido personal de los equipos de los
clientes (archivos, fotos, correos, conversaciones); tampoco datos bancarios o financieros, ni
informaci\'on de terceros que permita identificarlos. Si durante la observaci\'on se llegara a ver
contenido personal de un equipo, no ser\'a capturado ni registrado.

\subsection*{3. Riesgos y c\'omo se reducen}
La participaci\'on no implica riesgos f\'isicos. Los riesgos posibles son sentirse inc\'omodo al
hablar de su trabajo o de problemas internos de la organizaci\'on, y que la informaci\'on se use de
forma indebida. Para reducirlos: puede omitir cualquier pregunta, pausar o terminar la actividad
cuando quiera; el equipo usar\'a el c\'odigo de participante indicado anteriormente en lugar de su
nombre; y ninguna opini\'on suya se comunicar\'a a terceros de forma que lo identifique.

\subsection*{4. Beneficios}
No recibir\'a pago ni beneficio econ\'omico. Su colaboraci\'on ayudar\'a a que la propuesta de
sistema responda a las necesidades reales de su trabajo y contribuir\'a a la formaci\'on de futuros
profesionales de la Ingenier\'ia de Software.

\subsection*{5. Confidencialidad y protecci\'on de sus datos}
Sus datos personales ser\'an tratados conforme a la Ley Org\'anica de Protecci\'on de Datos
Personales del Ecuador (LOPDP), que le reconoce, entre otros, los derechos de acceso,
rectificaci\'on, eliminaci\'on y oposici\'on sobre sus datos.
\begin{itemize}[nosep]
  \item \textbf{Lo que nunca se publica:} su nombre, n\'umero de c\'edula, firma, rostro, voz,
        contacto ni su cargo espec\'ifico. En todo documento, presentaci\'on o repositorio p\'ublico
        usted aparecer\'a solo con su c\'odigo de participante o con un rol gen\'erico (por ejemplo,
        ``Responsable del taller'').
  \item \textbf{D\'onde se guarda lo identificable:} este formulario firmado y las grabaciones
        originales se guardan en una carpeta institucional privada (cuenta \texttt{@uteq.edu.ec}),
        compartida solo con el docente responsable. En el repositorio p\'ublico del equipo
        (ISR401-ContX) no se sube nada que identifique a nadie: solo material anonimizado, el
        manifiesto de las evidencias privadas y una copia de este formulario con nombre, c\'edula,
        firma y n\'umero de tel\'efono cubiertos.
  \item \textbf{Cu\'anto tiempo:} el material identificable se conservar\'a como m\'aximo 24 meses
        desde el cierre del per\'iodo acad\'emico y luego ser\'a destruido. Solo se conservar\'an los
        instrumentos en blanco y los resultados totalmente anonimizados.
\end{itemize}

\subsection*{6. Participaci\'on voluntaria y derecho a retirarse}
Su participaci\'on es libre. Puede negarse a responder cualquier pregunta y retirarse en cualquier
momento, sin dar explicaciones y sin ninguna consecuencia laboral, acad\'emica ni de otro tipo. Si
se retira, puede pedir que se eliminen sus datos usando el desprendible de la \'ultima p\'agina o
escribiendo a los contactos indicados al inicio; se eliminar\'an siempre que a\'un no se hayan
integrado en resultados agregados o anonimizados de forma irreversible. Este formulario est\'a
dirigido a personas de 18 a\~nos o m\'as; las personas menores de edad no participan en el estudio.

\subsection*{7. Autorizaciones espec\'ificas}
Marque S\'i o No en cada fila. Puede aceptar unas y rechazar otras; solo la primera es necesaria
para participar.

\begin{tabularx}{\linewidth}{@{}Ycc@{}}
  \toprule
  \textbf{Autorizaci\'on} & \textbf{S\'i} & \textbf{No} \\
  \midrule
  a) Acepto participar en la actividad que se me ha explicado. & $\square$ & $\square$ \\
  b) Autorizo la grabaci\'on en audio de mi participaci\'on. & $\square$ & $\square$ \\
  c) Autorizo la grabaci\'on en video de mi participaci\'on. & $\square$ & $\square$ \\
  d) Autorizo fotograf\'ias fechadas del entorno de trabajo, sin mi rostro, sin documentos
     legibles, sin pantallas que muestren datos de clientes y sin ubicaci\'on GPS. & $\square$ & $\square$ \\
  e) Autorizo que se citen textualmente frases m\'ias, sin datos que me identifiquen. & $\square$ & $\square$ \\
  f) Autorizo el uso de la informaci\'on anonimizada en los trabajos acad\'emicos de la asignatura. & $\square$ & $\square$ \\
  \bottomrule
\end{tabularx}

\smallskip
Ninguna de estas autorizaciones permite publicar datos que lo identifiquen.

\subsection*{8. Declaraci\'on y firmas}
Declaro que tengo 18 a\~nos o m\'as; que he le\'ido este formulario, o me lo han le\'ido; que pude
hacer preguntas y recib\'i respuestas claras; y que comprendo el objetivo, las actividades, los
riesgos, los beneficios, el manejo de mis datos y mi derecho a retirarme. Acepto participar de forma
voluntaria seg\'un las opciones que marqu\'e en la Secci\'on 7. Se firman dos ejemplares: uno queda en
mi poder y otro en el expediente restringido del proyecto.

\vspace{1.2cm}
\noindent
\begin{tabular}{@{}p{7.4cm}@{\hspace{0.8cm}}p{7.4cm}@{}}
  \rule{7.2cm}{0.4pt} & \rule{7.2cm}{0.4pt} \\
  \centering Nombres y apellidos del participante & \centering\arraybackslash Firma del participante \\[1cm]
  \rule{7.2cm}{0.4pt} & \rule{7.2cm}{0.4pt} \\
  \centering N\'umero de c\'edula & \centering\arraybackslash Lugar y fecha (dd/mm/aaaa) \\[1cm]
  \rule{7.2cm}{0.4pt} & \rule{7.2cm}{0.4pt} \\
  \centering Nombre y firma del estudiante que aplica el instrumento &
  \centering\arraybackslash Nombre y firma del testigo (opcional; obligatorio si el formulario se ley\'o en voz alta) \\
\end{tabular}

\subsection*{9. Revocaci\'on del consentimiento (opcional) --- desprendible para el participante}
Conserve esta hoja. Si m\'as adelante desea retirarse del estudio, compl\'etela y entr\'eguela a
cualquier integrante del equipo o env\'ie su contenido por correo a los contactos indicados en la
primera p\'agina.

\medskip
\noindent Yo, \rule{6cm}{0.4pt}, participante con c\'odigo \rule{2.5cm}{0.4pt}, retiro mi
consentimiento para participar en el proyecto Sistema de gesti\'on de reparaci\'on de equipos de
c\'omputo y solicito:
\begin{itemize}[nosep]
  \item[$\square$] Eliminar mis datos identificables (formulario firmado, grabaciones, fotograf\'ias).
  \item[$\square$] Eliminar tambi\'en mis respuestas anonimizadas, siempre que a\'un no se hayan
                   integrado en resultados agregados.
\end{itemize}

\vspace{1.2cm}
\noindent
\begin{tabular}{@{}p{7.4cm}@{\hspace{0.8cm}}p{7.4cm}@{}}
  \rule{7.2cm}{0.4pt} & \rule{7.2cm}{0.4pt} \\
  \centering Firma & \centering\arraybackslash Fecha (dd/mm/aaaa) \\
\end{tabular}


\end{document}
